\section{Pendahuluan}

Perangkat medis \textit{wearable} memungkinkan pemantauan elektrokardiogram (EKG) secara waktu nyata, tetapi kanal akuisisinya harus mempertahankan morfologi sinyal pada daya terbatas. Rantai akuisisi umumnya memuat penguat awal, filter analog, dan konverter analog-ke-digital \cite{Bio_signal_Processing,signal_processing_block}. Filter diperlukan karena gerakan tubuh, aktivitas di sekitar elektrode, dan interferensi jala-jala dapat mendistorsi bentuk gelombang EKG serta menurunkan ketepatan ekstraksi fitur. Kebutuhan ini menjadi menantang karena pita EKG berada pada frekuensi rendah, sekitar 0,5--250 Hz, sehingga implementasi filter terintegrasi harus menghasilkan frekuensi karakteristik rendah tanpa kapasitor yang berlebihan dan tanpa konsumsi daya tinggi \cite{Advanced_low_power_filter,Noise_ECG}.

Filter OTA-C/Gm-C sesuai untuk kebutuhan tersebut karena frekuensi karakteristiknya dapat ditentukan oleh \textit{transconductance} dan kapasitansi serta ditata ulang secara elektronik melalui arus bias \cite{OTA-C,gm_c}. Operasi MOS pada inversi lemah dapat menekan arus, tetapi tegangan ambang membatasi \textit{headroom} pada suplai rendah. Masukan melalui terminal \textit{bulk} mengurangi kendala tersebut dan telah digunakan pada filter bertegangan 0,3--0,5 V \cite{BD_4LPF,BulkInputOTA}. Penelitian terdahulu juga menerapkan \textit{bulk-driven multiple-input} OTA-C pada LPF Chebyshev orde lima untuk mereduksi artefak gerak EKG \cite{MIOTA-OTA}.

Kontribusi penelitian ini adalah perancangan dan verifikasi \textit{post-layout} sebuah \textit{biquad low-pass filter} orde dua berbasis \textit{bulk-driven} OTA-C pada CMOS 130 nm untuk pemrosesan EKG. Evaluasi mencakup karakterisasi rangkaian, ekstraksi parasitik, variasi PVT, serta pengujian EKG pada domain waktu dan frekuensi. Fokusnya ialah hubungan arus bias skala nA dengan frekuensi \textit{cut-off}, konsumsi daya, kualitas sinyal, dan keterbatasan operasi \textit{subthreshold}.
